\subsection{Exterior multiplication}

Let E be a hilbertian space. For every integer \(n \geqslant 0\) , let \(\theta_{n}\) be the canonical mapping from \(\mathbf{T}^{n}(\mathbf{E})\) onto \(\symbf{\Lambda}^{n}(\mathbf{E})\) ; then

\[
\theta_ {n} (x _ {1} \otimes \ldots \otimes x _ {n}) = x _ {1} \wedge \ldots \wedge x _ {n}\tag{31}
\]

for \(x_{1}, \ldots, x_{n}\) in E. Let \(p\) and \(q\) be two positive integers; on account of formulas (30) and (31) we have

\[
u \wedge v = \theta_ {p + q} \left(\frac {1}{(p !) ^ {1 / 2}} \psi_ {p} (u) \otimes \frac {1}{(q !) ^ {1 / 2}} \psi_ {q} (v)\right)\tag{32}
\]

for \(u \in \Lambda^p(\mathbf{E})\) and \(v \in \Lambda^q(\mathbf{E})\) . Since \(\| \theta_n \| \leqslant (n!)^{1/2}\) , we get the inequality

\[
\| u \wedge v \| \leqslant \left(\frac {(p + q) !}{p ! q !}\right) ^ {1 / 2} \| u \| \cdot \| v \|\tag{33}
\]

for \(u \in \symbf{\Lambda}^p(\mathrm{E})\) and \(v \in \symbf{\Lambda}^q(\mathrm{E})\) . Consequently, the mapping \((u, v) \mapsto u \wedge v\) extends by continuity to a bilinear mapping from \(\hat{\symbf{\Lambda}}^p(\mathrm{E}) \times \hat{\symbf{\Lambda}}^q(\mathrm{E})\) into \(\hat{\symbf{\Lambda}}^{p + q}(\mathrm{E})\) , with a norm at most equal to \(\left(\frac{(p + q)!}{p!q!}\right)^{1/2}\) (cf.~V, p.~73, exerc. 2). We again denote this by \((u, v) \mapsto u \wedge v\) .

PROPOSITION 6. --- Let E be a hilbertian space. We have

\[
\| x \wedge u \| \leqslant \| x \| \cdot \| u \|\tag{34}
\]

for \(x\in \mathbf{E}\) and \(u\in \hat{\Lambda} (\mathrm{E})\)

It is clearly enough to consider the case \(\| x\| = 1\)

Let F be the hilbertian subspace of E consisting of all vectors orthogonal to x. Since E is the hilbertian sum of F and the line K.x, it follows from the corollary of V, p.~34 that the mapping \((v,w)\mapsto v+x\wedge w\) is a hilbertian space isomorphism from \(\hat{\mathbf{A}}(\mathbf{F})\oplus\hat{\mathbf{A}}(\mathbf{F})\) onto \(\hat{\mathbf{A}}(\mathbf{E})\) . If \(u=v+x\wedge w\) with v,w in \(\hat{\mathbf{A}}(\mathbf{F})\) , we have \(x\wedge u=x\wedge v\) , hence \(\|x\wedge u\|=\|v\|\leqslant(\|v\|^{2}+\|w\|^{2})^{1/2}=\|u\|\) .

COROLLARY 1. --- a) Let \(x_1, \ldots, x_n\) be elements of the hilbertian space E. We have

\[
\left\| x _ {1} \wedge \dots \wedge x _ {n} \right\| \leqslant \left\| x _ {1} \right\| \dots \left\| x _ {n} \right\|,\tag{35}
\]

where equality holds only if one of the \(x_{i}\) is null, or if the sequence \((x_{1},\dots,x_{n})\) is orthogonal.

\begin{enumerate}
\def\labelenumi{\alph{enumi})}
\setcounter{enumi}{1}
\tightlist
\item
  Let \(x_{1}, \ldots, x_{n}, y_{1}, \ldots, y_{n}\) be elements of the hilbertian space E. We have
\end{enumerate}

\[
\left| \det (\langle x _ {i}, y _ {j} \rangle) \right| \leqslant \| x _ {1} \| \dots \| x _ {n} \| \cdot \| y _ {1} \| \dots \| y _ {n} \|;\tag{36}
\]

if the vectors \(x_{i}\) and \(y_{i}\) are non-null; equality holds in (36) if and only if \((x_{1},\ldots ,x_{n})\) and \((y_{1},\dots,y_{n})\) are each an orthogonal basis of the same vector subspace of \(\mathbf{E}\) .

The inequality (35) follows from prop. 6 by induction on n; the inequality (36) can be deduced by applying the Cauchy-Schwarz inequality in \(\hat{\Lambda}^{n}(E)\) and the formula (26) of V, p.~33.

Suppose that the sequence \((x_{1},\dots,x_{n})\) is orthogonal. Then

\[
\left\| x _ {1} \wedge \dots \wedge x _ {n} \right\| ^ {2} = \det (\langle x _ {i} | x _ {j} \rangle) = \prod_ {i = 1} ^ {n} \left\| x _ {i} \right\| ^ {2}
\]

since \(\langle x_i|x_j\rangle = 0\) for \(i\neq j\) .

Now suppose that the vectors \(x_{1}, \ldots, x_{n}\) are not null and do not form an orthogonal sequence. Since \(\| x_{1} \wedge \ldots \wedge x_{n} \|\) depends symmetrically on the vectors \(x_{1}, \ldots, x_{n}\) , we can assume that \(x_{1}\) is not orthogonal to the subspace \(F\) of \(E\) generated by \(x_{2}, \ldots, x_{n}\) , and that \(F\) is not simply 0. We can decompose \(x_{1}\) in the form \(x_{1}' + y\) with \(y \neq 0\) in \(F\) and \(x_{1}'\) orthogonal to \(F\) ; then \(\| x_{1}' \| < \| x_{1} \|\) . But

\[
x _ {1} \wedge x _ {2} \wedge \ldots \wedge x _ {n} = x _ {1} ^ {\prime} \wedge x _ {2} \wedge \ldots \wedge x _ {n},
\]

and so

\[
\begin{array}{r l} \left\| x _ {1} \wedge \dots \wedge x _ {n} \right\| & \leqslant \left\| x _ {1} ^ {\prime} \right\| \left\| x _ {2} \right\| \dots \left\| x _ {n} \right\| \\ & <   \left\| x _ {1} \right\| \left\| x _ {2} \right\| \dots \left\| x _ {n} \right\|. \end{array}
\]

Suppose that the vectors \(x_{i}\) and the vectors \(y_{i}\) are not null. The equality in relation (36) is equivalent to the conjunction of the equalities

\[
\left| \left\langle x _ {1} \wedge \dots \wedge x _ {n} \mid y _ {1} \wedge \dots \wedge y _ {n} \right\rangle \right| = \| x _ {1} \wedge \dots \wedge x _ {n} \| \cdot \| y _ {1} \wedge \dots \wedge y _ {n} \|\tag{37}
\]

\[
\left\| x _ {1} \wedge \dots \wedge x _ {n} \right\| = \left\| x _ {1} \right\| \dots \left\| x _ {n} \right\|, \quad \left\| y _ {1} \wedge \dots \wedge y _ {n} \right\| = \left\| y _ {1} \right\| \dots \left\| y _ {n} \right\|.\tag{38}
\]

By the first part of the proof, the equalities (38) imply that each of the sequences \((x_{1},\ldots ,x_{n})\) and \((y_{1},\ldots ,y_{n})\) is orthogonal, which in turn implies that \(x_{1}\wedge \ldots \wedge x_{n}\neq 0\) and that \(y_{1}\wedge \ldots \wedge y_{n}\neq 0\) . Under these conditions, relation (37) implies that there exists a scalar \(\lambda \neq 0\) such that \(y_{1}\wedge \ldots \wedge y_{n} = \lambda x_{1}\wedge \ldots \wedge x_{n}\) (V, p.~3, Remark 1); in other words, that \((x_{1},\dots,x_{n})\) and \((y_{1},\dots,y_{n})\) are bases of the same vector subspace of E (A, III, § 11, No.~13).

COROLLARY 2. --- Let \((a_{ij})_{1\leqslant i,j\leqslant n}\) be a hermitian matrix, with complex elements, and with determinant D. Suppose that the inequality

\[
\sum_ {i, j = 1} ^ {n} a _ {i j} \overline {{z}} _ {i} z _ {j} \geqslant 0\tag{39}
\]

holds for all complex numbers \(z_{1}, \ldots, z_{n}\) . Then

\[
0 \leqslant \mathrm{D} \leqslant a _ {1 1} \dots a _ {n n}.\tag{40}
\]

Suppose \(\mathbf{D}\) is non-null; the equality \(\mathbf{D} = a_{11} \ldots a_{nn}\) holds if and only if \(a_{ij} = 0\) for all \(i \neq j\) .

Let \(\Phi\) be the hermitian form on the vector space \(\mathbf{C}^n\) given by

\[
\Phi (\mathbf {z}, \mathbf {z} ^ {\prime}) = \sum_ {i, j = 1} ^ {n} a _ {i j} \overline {{{{z}}}} _ {i} z _ {j} ^ {\prime}
\]

for \(\mathbf{z} = (z_1, \dots, z_n)\) and \(\mathbf{z}' = (z_1', \dots, z_n')\) in \(\mathbf{C}^n\) . By hypothesis, \(\Phi\) is positive.

First assume that \(\Phi\) is separating, that is, that D is non-null. If \((e_1, \ldots, e_n)\) is the canonical basis of \(\mathbf{C}^n\) , we have \(\Phi(\mathbf{e}_i, \mathbf{e}_j) = a_{ij}\) , and cor. 2 follows immediately from cor. 1, \(a)\) by putting \(x_i = \mathbf{e}_i\) .

Since \(a_{ii} = \Phi(e_i, e_i) \geqslant 0\) , we also have inequality (40) if \(D = 0\) .

COROLLARY 3 (« Hadamard's inequalities »). --- Let \((a_{ij})_{1 \leqslant i,j \leqslant n}\) be a matrix with complex elements, and with determinant D. Put

\[
c _ {i} = \left(\sum_ {j = 1} ^ {n} | a _ {i j} | ^ {2}\right) ^ {1 / 2} \quad \text {for} \quad 1 \leqslant i \leqslant n,
\]

and \(m = \sup_{i,j}|a_{ij}|\) . Then we have

\[
| \mathbf {D} | \leqslant c _ {1} \dots c _ {n} \leqslant m ^ {n}. n ^ {n / 2}.\tag{41}
\]

If \(\mathbf{D} \neq 0\) , in order that \(|\mathbf{D}| = c_1 \ldots c_n\) it is necessary and sufficient that the rows \(y_i = (a_{ij})_{1 \leqslant j \leqslant n}\) of the matrix \((a_{ij})_{1 \leqslant i,j \leqslant n}\) are two by two orthogonal vectors.

Let the space \(C^{n}\) be assigned a scalar product defined by

\[
\langle \mathbf {z} | \mathbf {z} ^ {\prime} \rangle = \sum_ {i = 1} ^ {n} \overline {{z}} _ {i} z _ {i} ^ {\prime}.
\]

Let \((\mathbf{x}_{1},\ldots,\mathbf{x}_{n})\) be the canonical basis of \(C^{n}\) and \(y_{i}\) the vector with components \(a_{ij}\) for \(1\leqslant j\leqslant n\) . We have \(\|x_{i}\|=1\) and \(\|y_{i}\|=c_{i}\) for \(1\leqslant i\leqslant n\) ; also \(\langle x_{i}|y_{j}\rangle=a_{ji}\) . The inequality \(|D|\leqslant c_{1}\ldots c_{n}\) and the condition of equality are then particular cases of V, p.~35, cor. 1. Obviously we have \(c_{i}\leqslant m.n^{1/2}\) , hence \(c_{1}\ldots c_{n}\leqslant m^{n}.n^{n/2}\) .

